\section{Trabajos Relacionados}
\label{sec:trabajos_relacionados}

La presente propuesta se fundamenta en la línea de investigación de la especificación formal con VDM++ y su aplicación a sistemas de implementación crítica y de gran escala.

Respecto al soporte de herramientas, Fitzgerald, Larsen y Sahara~\cite{fitzgerald2008} describen VDMTools, un conjunto de herramientas que analizan modelos expresados en VDM y sus dialectos VDM++ y VDM-RT, y que ofrecen verificación estática, generación automática de código, mapeo a UML, análisis de cobertura de pruebas y depuración. Este trabajo destaca que VDMTools está orientado al desarrollo rentable de modelos formales en entornos industriales, lo que sustenta la elección de esta herramienta para validar el algoritmo propuesto.

En el ámbito de los sistemas críticos, Verhoef, Larsen y Hooman~\cite{verhoef2006} extienden VDM++ para modelar y validar sistemas embebidos distribuidos en tiempo real, definiendo una semántica operacional constructiva validada con VDMTools. De forma complementaria, Verhoef~\cite{verhoef2005} analiza el uso de VDM++ para especificar sistemas de tiempo real, argumentando que los métodos formales deben integrarse a los flujos de trabajo industriales y permitir el análisis temprano de propiedades. Ambos trabajos justifican la aplicación de VDM++ a un componente cuya correctitud es crítica, como el algoritmo de planificación de viajes propuesto.

Finalmente, Mochio y Araki~\cite{mochio2011} vinculan VDM++ con el desarrollo formal ágil y escalable de software de misión crítica, incluyendo la descripción de arquitecturas y la verificación de especificaciones. Este enfoque respalda el requisito de escalabilidad de la solución, en la medida en que el modelo formal puede acompañar la evolución del sistema.

A diferencia de los trabajos citados, centrados en sistemas de tiempo real o en la infraestructura de herramientas, esta propuesta aplica VDM++ a la verificación de un algoritmo de planificación de rutas de transporte público, en el contexto específico de la ciudad de Arequipa.
